\documentclass[12pt]{article}
\usepackage[T1]{fontenc}
\usepackage[utf8]{inputenc}
\usepackage{lmodern}
\usepackage{textcomp}
\usepackage{enumitem}
\usepackage{geometry}
\usepackage{graphicx}
\usepackage{amsmath, amssymb}
\geometry{margin=1in}
\begin{document}
\begin{center}
History - Undergraduate Level Assessment 20 \
Total Marks: 40 \
Answer all questions.
\end{center}

\section*{Multiple Choice (10 marks)}
\begin{enumerate}[label=\arabic*.]
\item What distinguishes the trade patterns of the Mesolithic from those of the Upper Paleolithic?
\begin{enumerate}[label=(\alph*)]
    \item in the Mesolithic, goods traveled less far and fewer goods were exchanged
    \item in the Mesolithic, goods traveled farther but the same number of goods were exchanged
    \item in the Mesolithic, goods traveled less far but the same number of goods were exchanged
    \item in the Mesolithic, goods traveled the same distance but more goods were exchanged
    \item in the Mesolithic, goods traveled farther but fewer goods were exchanged
\end{enumerate}
\item Which of the following characteristics of modern humans is the most ancient, pre-dating the Late Stone Age, the Upper Paleolithic and even anatomically modern humans?
\begin{enumerate}[label=(\alph*)]
    \item establishment of permanent settlements
    \item development of agriculture
    \item domestication of animals
    \item creation of written language
    \item symbolic expression through the production of art
\end{enumerate}
\item The \_\_\_\_\_\_\_ developed one of the earliest kingdoms in South America, reaching its height at around \_\_\_\_\_\_\_\_\_.
\begin{enumerate}[label=(\alph*)]
    \item Inca; AD 200
    \item Moche; AD 400
    \item Olmec; AD 100
    \item Maya; AD 400
    \item Zapotec; AD 500
\end{enumerate}
\item When did the first pharaohs emerge in Egypt?
\begin{enumerate}[label=(\alph*)]
    \item 3100 B.P.
    \item 4100 B.P.
    \item 5100 B.P.
    \item 6100 B.P.
\end{enumerate}
\item The commoners of the rural Aztec villages of Capilco and Cuexcomate:
\begin{enumerate}[label=(\alph*)]
    \item had to pay a high tax in goods and labor to the local governor.
    \item lived under constant threat of being sold into slavery.
    \item were mostly left alone and did fairly well for themselves.
    \item were under the direct control of the high priest, who demanded regular sacrifices.
    \item were heavily policed by the king's guards to prevent any form of uprising.
\end{enumerate}
\end{enumerate}

\section*{True / False (10 marks)}
\textit{Answer True or False}
\begin{enumerate}[label=\arabic*.]
\setcounter{enumi}{5}
\item True or False: Even though the book had barely hinted at human evolution, it quickly became central to the debate as mental and moral qualities were seen as spiritual aspects of the immaterial soul, and it was believed that animals did not have spiritual qualities. This conflict could be reconciled by supposing there was some supernatural intervention on the path leading to humans, or viewing evolution as a purposeful and progressive ascent to mankind's position at the head of nature.
\item True or False: This approach did not serve France well in the war, as the colonies were indeed lost, but although much of the European war went well, by its end a different answer
\item True or False: Hans Bielenstein writes that as far back as the Han dynasty (202 BCE-220 CE), the Han Chinese government "maintained the fiction" that the a different answer administering the various "Dependent States" and oasis city-states of the Western Regions (composed of the Tarim Basin and oasis of Turpan) were true Han representatives due to the Han government's conferral of Chinese seals and seal cords to them.
\item True or False: Because of a low literacy rate among the public (about 2-3\% until the early 19 th century and just about 15\% at the end of 19 th century),[citation needed] ordinary a different answer as "special request-writers" (arzuhâlcis) to be able to communicate with the government.
\item True or False: Earth was initially molten due to extreme volcanism and frequent collisions with other bodies.
\end{enumerate}

\section*{Long Form Response (20 marks)}
\textit{Answer each question with a clear and complete explanation.}
\begin{enumerate}[label=\arabic*.]
\setcounter{enumi}{10}
\item Read the following passage and answer the question: Several films have featured their songs performed by other artists. A version of "Somebody to Love" by Anne Hathaway was in the 2004 film Ella Enchanted. In 2006, Brittany Murphy also recorded a cover of the same song for the 2006 film Happy Feet. In 2001, a version of "The Show Must Go On" was performed by Jim Broadbent and Nicole Kidman in the film musical Moulin Rouge!. The 2001 film A Knight's Tale has a version of "We Are the Champions" performed by Robbie Williams and Queen; the film also features "We Will Rock You" played by the medieval audience. Question: ho sang a version of Queen's Somebody to Love in 2004's Ella Enchanted?
\item Read the following passage and answer the question: The Space Age is a period encompassing the activities related to the Space Race, space exploration, space technology, and the cultural developments influenced by these events. The Space Age began with the development of several technologies that culminated with the launch of Sputnik 1 by the Soviet Union. This was the world's first artificial satellite, orbiting the Earth in 98.1 minutes and weighing in at 83 kg. The launch of Sputnik 1 ushered a new era of political, scientific and technological achievements that became known as the Space Age. The Space Age was characterized by rapid development of new technology in a close race mostly between the United States and the Soviet Union. The Space Age brought the first human spaceflight during the Vostok programme and reached its peak with the... Question: How many people watched the Apollo 11 landing?
\end{enumerate}

\vfill
\noindent\textit{End of Paper}
\end{document}